\chapter{Gaussian, fermionic, and matchgate compression}
\label{ch:gaussian}
\index{Gaussian circuit}
\index{matchgate}

\section{Polynomial covariance representations}
Gaussian bosonic and fermionic states are specified by first and second moments rather than by all Fock-basis amplitudes. A bosonic covariance matrix has size polynomial in the number of modes, and a fermionic Gaussian state is represented by an antisymmetric Majorana covariance matrix. Gaussian unitaries act by symplectic or orthogonal transformations:
\[
\Gamma\longmapsto S\Gamma S^T
\quad\text{or}\quad
\Gamma\longmapsto O\Gamma O^T.
\]
This is a compact internal encoding even when the corresponding Hilbert-space packet is exponentially large.

\section{Hybrid symmetry and Gaussianity}
Number conservation, parity, translation symmetry, and point-group symmetry can further block the covariance matrices. A simulator can therefore combine two reductions:
\[
\text{symmetry packetization}
+
\text{Gaussian internal encoding}.
\]
The first removes unreachable sectors. The second compresses states within the remaining sectors.

\section{Matchgate circuits}
Nearest-neighbor matchgate circuits have polynomial classical representations through fermionic Gaussian dynamics. Geometry matters: adding SWAP, fermionic SWAP, long-range gates, or a non-matchgate resource changes the simulation problem. A symmetry-collapse atlas should therefore record the gate library and connectivity as part of the model, not as implementation details.

\section{Low-rank non-Gaussian corrections}
One research direction is to represent a weakly non-Gaussian state as
\[
\rho
=
\sum_{a=1}^R c_a\rho_a^{\rm Gauss}
\]
with small Gaussian rank $R$, or as a Gaussian reference plus sparse packet corrections. A polynomial collapse requires $R$ and the update cost to remain polynomial.

\begin{frontierbox}
Gaussianity is an internal representation class, not a symmetry projector. A complete hybrid theorem must prove closure of the chosen compressed class under every allowed gate and bound the growth of any non-Gaussian residue.
\end{frontierbox}

\section{Few-particle versus Gaussian routes}
Fixed-particle simulation and Gaussian simulation are independent. A fixed two-fermion state may be non-Gaussian but live in a polynomial-dimensional sector. A half-filled Gaussian state may live in an exponential number sector but have a polynomial covariance description. The two-axis model in \cref{ch:compressed-packets} is designed to capture this distinction.
